\documentclass[10pt,a4paper]{amsart}
\usepackage[margin=19mm]{geometry}
\usepackage{amsmath,amssymb,mathtools}
\usepackage[scr=rsfs,cal=euler]{mathalfa}
\usepackage{xfrac}
\usepackage[unicode,hidelinks,bookmarks=false]{hyperref}
\newcommand{\ie}{i.e.\ }
\newcommand{\cf}{cf.\ }
\newcommand{\resp}{respectively\ }
\newcommand{\Subsection}{\subsection}
\providecommand{\nobreakdash}{\nobreak\hbox{-}}
\setlength{\emergencystretch}{2em}
\multlinegap0pt
\usepackage{tikz}
\usetikzlibrary{cd}
\usepackage{amssymb}
\newtheorem{prop}{Proposition}
\title{On a Completion of Cohomological Functors Generalising Tate Cohomology I}
\author{Max Gheorghiu}
\date{Source: arXiv:2405.03610v3 (CC BY 4.0)}
\begin{document}
\maketitle

\noindent\emph{Source and licence.} On a Completion of Cohomological Functors Generalising Tate Cohomology I. Version arXiv:2405.03610v3 (CC BY 4.0), distributed by arXiv under the Creative Commons Attribution 4.0 International licence, http://creativecommons.org/licenses/by/4.0/, as recorded in the arXiv licence metadata for this exact version and shown on the abstract page https://arxiv.org/abs/2405.03610v3. Full licence notice: \texttt{licenses/arxiv-2405.03610-CC-BY-4.0.txt}.\par\smallskip

\noindent\textit{Editorial scope.} Counted unit: the article's prop prop:transitioning (source0.tex lines 847--849). The article's complete original proof of it is delivered with it (source0.tex lines 851--871, one \texttt{proof} environment of 21 lines). No mathematics is written, repaired or re-derived: the statement and the proof are the article's own lines, in the article's own notation, and no retained line is substituted or re-flowed.  Retained uncounted as the source-local context the proof needs: lines 556--561; lines 841--843.  Nothing was omitted from any retained span, so the package carries no omission record.  No retained line contains a citation, so the wrapper carries no bibliography and no bibliography entry.  The retained lines contain the article's own diagram code, which the wrapper typesets with the standard tikz packages; no image file is delivered and the package declares no asset dependency.  Editorial formatting only, in addition: an AMS \texttt{amsart} preamble that restates the article's own declarations and macros used by the retained lines, namely the environments \texttt{prop} on separate counters, together with the standard packages those lines need.  The retained environments print in this document as Proposition 1 in order of appearance, whereas the article prints the same items with the numbering of its own full document.  The complete native source of the article as distributed in the arXiv e-print for this exact version is bundled unchanged as \texttt{sources/158731/source0.tex}.
\begin{equation}\label{diag:complexlift}
\begin{tikzcd}
    0 \arrow[r] & \widetilde{M}_{k+1} \arrow[r] \arrow[d, "\widetilde{f}_{k+1}"] & M_k \arrow[r] \arrow[d, "f_k"] & \widetilde{M}_k \arrow[r] \arrow[d, "\widetilde{f}_k"] & 0 \\
    0 \arrow[r] & \widetilde{N}_{k+1} \arrow[r] & N_k \arrow[r] & \widetilde{N}_k \arrow[r] & 0
\end{tikzcd}
\end{equation}

\begin{prop}\label{prop:biadditive}
The bifunctor $\circ: \mathrm{Hom}_{\mathcal{C}}(B, C) \times \mathrm{Hom}_{\mathcal{C}}(A, B) \rightarrow \mathrm{Hom}_{\mathcal{C}}(A, C)$ descends to a bifunctor $\circ: [B, C]_{\mathcal{C}} \times [A, B]_{\mathcal{C}} \rightarrow [A, C]_{\mathcal{C}}$ that is associative and additive in both variables.
\end{prop}

\begin{prop}\label{prop:transitioning}
There exists a homomorphism $t_{A, B}: [A, B]_{\mathcal{C}} \rightarrow [\widetilde{A}_1, \widetilde{B}_1]_{\mathcal{C}}$ with the property that $t_{B, C}(-) \circ t_{A, B}(-) = t_{A, C}(- \circ -)$.
\end{prop}

\begin{proof}
Consider the following version of Diagram~\ref{diag:complexlift}
\begin{equation}\label{diag:simplelift}
\begin{tikzcd}
    0 \arrow[r] & \widetilde{A}_1 \arrow[r, "\iota_A"] \arrow[d, "\widetilde{f}_1"] & A_0 \arrow[r, "\pi_A"] \arrow[d, "f_0"] & A \arrow[r] \arrow[d, "f"] & 0 \\
    0 \arrow[r] & \widetilde{B}_1 \arrow[r, "\iota_B"] & B_0 \arrow[r, "\pi_B"] & B \arrow[r] & 0
\end{tikzcd}
\end{equation}
where $A_0$ and $B_0$ are assumed to be projective. If $f_0': A_0 \rightarrow B_0$ and $\widetilde{f}_1': \widetilde{A}_1 \rightarrow \widetilde{B}_1$ are different lifts, then there exists a morphism $e: A_0 \rightarrow \widetilde{B}_1$ such that $\widetilde{f}_1-\widetilde{f}_1' = e \circ \iota_A$. As $A_0$ is projective, $\widetilde{f}_1$ and $\widetilde{f}_1'$ agree in $[\widetilde{A}_1, \widetilde{B}_1]_{\mathcal{C}}$. Therefore, there is a well defined homomorphism
\[ s_{A, B}: \mathrm{Hom}_{\mathcal{C}}(A, B) \rightarrow [\widetilde{A}_1, \widetilde{B}_1]_{\mathcal{C}}, f \mapsto \widetilde{f_1}+\mathcal{P}_{\mathcal{C}}(\widetilde{A}_1, \widetilde{B}_1) \, . \]
If $h \in \mathrm{Hom}_{\mathcal{C}}(B, C)$, then we see that $s_{B, C}(h) \circ s_{A, B}(f) = s_{A, C}(h \circ f)$. In particular, if $f \in \mathcal{P}_{\mathcal{C}}(A, B)$, then $s_{A, B}(f)$ = 0. Thus, $s_{A, B}$ descends to the desired homomorphism $t_{A, B}: [A, B]_{\mathcal{C}} \rightarrow [\widetilde{A}_1, \widetilde{B}_1]_{\mathcal{C}}$. \\

Because $s_{B, C}(-) \circ s_{A, B}(-) = s_{A, C}(- \circ -)$, the respective maps factor through quotient homomorphisms as in the commutative diagram
\begin{center}
\begin{tikzcd}[column sep = scriptsize]
	{\scriptstyle \mathrm{Hom}_{\mathcal{C}}(B{,} \, C) \times \mathrm{Hom}_{\mathcal{C}}(A{,} \, B)} \arrow[r] \arrow[rrrr, bend right = 10, "{\scriptstyle s_{B, C} \times s_{A, B}}"] \arrow[ddd, "{\scriptstyle - \circ -}"] & {\scriptstyle \mathrm{Hom}_{\mathcal{C}}(B{,} \, C) \times {[}A{,} \, B{]}_{\mathcal{C}}} \arrow[r] \arrow[rrr, bend left = 15, "{\scriptstyle s_{B, C} \times t_{A, B}}"] & {\scriptstyle {[}B{,} \, C{]}_{\mathcal{C}} \times {[}A{,} \, B{]}_{\mathcal{C}}} \arrow[rr, "{\scriptscriptstyle t_{B, C} \times t_{A, B}}" near start] & & {\scriptstyle {[}\widetilde{B}_1{,} \, \widetilde{C}_1{]} \times {[}\widetilde{A}_1{,} \, \widetilde{B}_1{]}} \arrow[ddd, "{\scriptstyle - \circ -}"] \\ \\ \\
	{\scriptstyle \mathrm{Hom}_{\mathcal{C}}(A{,} \, C)} \arrow[rr] \arrow[rrrr, bend right = 10, "s_{A, C}"] & & {\scriptstyle {[}A{,} \, C{]}_{\mathcal{C}}} \arrow[rr, "{\scriptstyle t_{A, C}}"] & & {\scriptstyle {[}\widetilde{A}_1{,} \, \widetilde{C}_1{]}}
\end{tikzcd}
\end{center}
Its right-hand side and Proposition~\ref{prop:biadditive} imply that $t_{B, C}(-) \circ t_{A, B}(-) = t_{A, C}(- \circ -)$.
\end{proof}

\end{document}
